% Owner: Aliza (dataset, annotations). Numbers come from results/report/numbers.tex.
\section{Dataset}
\label{sec:data}

\paragraph{Images.}
We collected \NumImages{} images in six categories where color carries meaning for people
with color vision deficiency (\cref{tab:dataset}): charts and graphs, maps, traffic signals
and road signs, ripe versus unripe food, color-coded items such as wires, bins, and LEDs, and
natural scenes. \NumWebImages{} images come from public web pages and \NumTeamPhotos{} are
photographs taken by the team. Every image is stored as an sRGB PNG whose longest side is at
most 1024~px; smaller source images are kept at their native size rather than upsampled.
Following the project's color conventions, images are decoded to linear RGB in $[0,1]$ before
any simulation, recoloring, or measurement, and every color difference is CIEDE2000
\cite{luo2001,sharma2005} computed in CIELAB under D65.

\paragraph{Splits.}
A tuning split of \NumTune{} images was drawn once, stratified by category with largest-remainder
allocation and a fixed seed. It is the only data used to tune the personalized method or to
make any modelling choice; the remaining \NumTest{} images form the test split on which every
number in \cref{sec:evaluation} is computed.

\begin{table}[t]
  \centering
  \small
  \begin{tabular}{@{}lrrrrr@{}}
\toprule
Category & Images & Tune & Test & Annot.\ test & Pairs (test) \\
\midrule
Charts & 40 & 6 & 34 & 9 & 26 \\
Maps & 30 & 5 & 25 & 9 & 27 \\
Traffic & 30 & 4 & 26 & 8 & 16 \\
Food & 40 & 6 & 34 & 9 & 22 \\
Coded items & 30 & 5 & 25 & 8 & 22 \\
Natural & 30 & 4 & 26 & 8 & 16 \\
\midrule
\textbf{Total} & 200 & 30 & 170 & 51 & 129 \\
\bottomrule
\end{tabular}

  \caption{Dataset composition. \emph{Annot.\ test}: test images with region-pair
  annotations. In addition, \NumAnnotTune{} tune images carry \NumPairsTune{} pairs, used
  only for tuning.}
  \label{tab:dataset}
\end{table}

\paragraph{Region-pair annotations.}
For an image to be useful to a viewer with CVD, specific regions must stay distinguishable:
the red and green bars of a chart, a lit signal lamp and an unlit one, a ripe and an unripe
fruit. We annotate such \emph{region pairs} $(A, B)$ as non-overlapping polygons in pixel
coordinates, each with a short note (\eg ``red pepper vs green pepper''). The test split has
\NumPairsTest{} pairs on \NumAnnotTest{} images, spread over all six categories, and the tune
split has \NumPairsTune{} pairs on \NumAnnotTune{} images.

Annotations were drafted with AI assistance and then reviewed by a team member. Candidate
regions came from graph-based segmentation~\cite{felzenszwalb2004}, each reduced to the largest
axis-aligned rectangle inside a single-color segment; thin structures (lines, small lamps)
used a seeded color flood fill, and scatter-plot clusters used a manual box. For each image a
person chose which two regions form a meaningful pair. Every pair was checked automatically
(regions inside the image, no overlap, and a normal-vision difference of at least
5~$\Delta E_{00}$) and visually on an overlay of the polygons. Because pairs were chosen to be
clearly different for normal vision, the dataset measures how much of that difference survives
CVD simulation, not near-threshold discrimination.

\paragraph{Metadata and licensing.}
The manifest records source, source URL, license, photographer, lighting, split, and whether
an image is annotated. Lighting was labelled for the team photographs (daylight, indoor, or
low light). Of the \NumWebImages{} web images, \NumOpenLicense{} have an open license (mostly Creative
Commons via Wikimedia Commons) or are public domain, and \NumNoLicense{} state no license. We therefore use the images for evaluation only and do
not redistribute them; the qualitative figure in this report uses a team photograph.
